% titel: Blatt 2, A8: Das Haber-Bosch-Verfahren
% skript: Mortimer Kap. (Katalyse, Ammoniaksynthese)

% Teil I: Werkzeuge, kein Lösungsweg.

Das Haber-Bosch-Verfahren ist die industrielle Synthese von Ammoniak aus den
Elementen. Wichtig sind die Reaktionsgleichung, die Identifikation des
begrenzenden Reaktanden und der Einsatzzweck.

\begin{werkzeugkasten}[Haber-Bosch-Verfahren]
Grundreaktion zwischen den Elementen (ausbalanciert, Katalysator + Energie):
\begin{equation}
  \mathrm{N_2} + 3\,\mathrm{H_2} \;\to\; 2\,\mathrm{NH_3}
\end{equation}
\begin{itemize}
  \item Die Gleichung ist nach dem Erhalt von N und H auszugleichen
        (Begriffe aus A7: Koeffizienten, begrenzender Reaktand).
  \item In der Technik liegt meist Stickstoff (aus der Luft, quasi
        unbegrenzt verfügbar) im Ueberschuss vor; der Katalysator wärmt die
        Umsetzung bei hohem Druck und Temperatur.
  \item Anwendung: Herstellung von Ammoniak, vor allem für Düngemittel.
  \item Begriffe ``Substrat'' und ``Edukt'': ``Edukt'' ist der
        Oberbegriff für alle Ausgangsstoffe einer Reaktion; ``Substrat'' ist
        bei katalytischen Reaktionen derjenige Stoff, der mithilfe des
        Katalysators umgewandelt wird.
\end{itemize}
\textit{Skript:} Mortimer Kap. (Katalyse, Ammoniaksynthese), Übung 2/8.
\end{werkzeugkasten}

\begin{achtung}[Darauf achten]
\begin{itemize}
  \item Gleichung stets ausbalancieren (Atomerhalt), bevor man Mengen nennt.
  \item Begrenzender Reaktand ist eine Sache der gegebenen Stoffmengen, nicht
        der Grösse der Koeffizienten allein.
  \item ``Substrat'' ist auf katalysierte Umsetzungen bezogen, ``Edukt'' ist
        der allgemeinere Begriff.
\end{itemize}
\end{achtung}
